\documentclass[11pt]{article}
\usepackage{amsmath,amssymb,amsthm,hyperref}
\newtheorem{theorem}{Theorem}
\newtheorem{lemma}[theorem]{Lemma}
\newtheorem{corollary}[theorem]{Corollary}
\begin{document}
\title{Polynomial-size rational monotone models for the local comparison tensor}
\author{Research note; complete proof candidate}
\date{30 September 2026}
\maketitle

\begin{theorem}\label{thm:main}
There is an absolute constant $C$ such that for every $m\geq1$ there are
an integer $K\leq Cm^2$ and rational numbers $p_{j,q}$, $1\leq j\leq m$,
$1\leq q\leq K$, satisfying
\[
 0<p_{j,1}<p_{j,2}<\cdots<p_{j,K}<1\qquad(j\in[m])
\]
and
\begin{equation}\label{eq:moments}
 \frac1K\sum_{q=1}^K\prod_{j\in S}p_{j,q}=\frac1{|S|+1}
 \qquad(S\subseteq[m]).
\end{equation}
Consequently the equal-weight product mixture with parameters $p_{j,q}$
has precisely the comparison tensor induced by a uniform random ranking:
\[
 \frac1K\sum_{q=1}^K\prod_{j\in S}p_{j,q}
                         \prod_{j\notin S}(1-p_{j,q})
   =\int_0^1t^{|S|}(1-t)^{m-|S|}\,dt.
\]
\end{theorem}

\paragraph{Scope.}
Thus rationality, positivity, monotonicity in the distinguished face, equal
face weights, and \emph{all} local comparison probabilities do not force
an exponential number of faces. The exponential rational scalar-quadrature
obstruction cannot be extended to these local hypotheses alone. A proof of
an exponential per-die lower bound must use additional conditions involving
the order among the other dice. This theorem does not produce a
permutation-fair family.

The only quantitative input is the classical real Chebyshev-type quadrature
bound: degree $d$ admits an equal-weight interval quadrature with $O(d^2)$
nodes. For example, this is a special case of Gilboa--Peled,
\emph{Chebyshev-type quadratures for doubling weights}, Theorem 1.1,
\href{https://arxiv.org/abs/1507.01505}{arXiv:1507.01505}.
The following lemma makes the necessary distinctness and interiority
explicit, without assuming these additional properties in that theorem.

\begin{lemma}\label{lem:regular}
For every $m\geq1$, there is a degree-$m$ real equal-weight quadrature
with $K$ nodes satisfying
\[
 m^2\leq K=O(m^2),\qquad 0<u_1<\cdots<u_K<1.
\]
\end{lemma}
\begin{proof}
Take a degree-$(2m+2)$ real equal-weight rule with $K_0=O(m^2)$ nodes,
allowing repetitions and endpoints. It has at least $m+2$ distinct support
points: otherwise the square of the polynomial vanishing on its support
would have degree at most $2m+2$, zero quadrature sum, and strictly positive
integral. At least $m$ of the support points therefore lie in $(0,1)$.

Repeat the entire rule $\lceil m^2/K_0\rceil$ times. Its new cardinality
$K$ is at least $m^2$ and less than $m^2+K_0=O(m^2)$ (with a harmless
nonstrict inequality when $m^2/K_0$ is integral). Select $m$ coordinates at
distinct interior support points as pivot coordinates. The Jacobian of the
first $m$ moment sums with respect to these pivots is an invertible scaled
Vandermonde matrix. The analytic implicit function theorem therefore
allows arbitrary sufficiently small perturbations of the other, free,
coordinates, while uniquely adjusting the pivot coordinates to keep the
first $m$ moments exact.

All pivot coordinates remain interior and mutually distinct. Move each free
endpoint slightly inward and choose the free perturbations generically to
avoid collisions. Here is a precise verification that this is possible.
Free--free collisions are proper affine hyperplanes. A pivot--free collision
that holds at the starting rule is a nonzero analytic equation: varying only
the colliding free coordinate, the derivative of its matching pivot is $-1$,
by the invertible Jacobian and the equality of their moment-derivative
columns. Their difference thus has derivative $-2$. Collisions not present
at the starting point remain absent in a sufficiently small neighborhood.
A finite union of proper real analytic zero sets cannot cover the open set
of allowed inward perturbations. Sorting the resulting nodes gives the lemma.
\end{proof}

\begin{proof}[Proof of Theorem~\ref{thm:main}]
Choose a rule $u=(u_q)_{q=1}^K$ from Lemma~\ref{lem:regular}.
Choose pairwise disjoint sets $Q_1,\ldots,Q_m\subseteq[K]$ of size $m$;
unused indices are permitted. Approximate $u$ arbitrarily closely by a
strictly increasing rational vector $t=(t_q)_{q=1}^K\in(0,1)^K$.
For $1\leq s\leq m$, put
\[
 \mu_s(t)=\frac1K\sum_{q=1}^Kt_q^s,
 \qquad a_s(t)=\frac{1/(s+1)-\mu_s(t)}s.
\]
For each $j$, solve the $m$-by-$m$ Vandermonde system
\begin{equation}\label{eq:correct}
 \frac1K\sum_{q\in Q_j}f_j(q)t_q^{s-1}=a_s(t)
          \qquad(1\leq s\leq m),
\end{equation}
and set $f_j(q)=0$ outside $Q_j$.
The solution is rational. Moreover it depends continuously on $t$ near $u$,
and it tends to zero as $t\to u$, because $a_s(u)=0$ and the limiting
Vandermonde matrices are invertible.
Set $p_{j,q}=t_q+f_j(q)$. For a sufficiently accurate rational approximation,
all columns remain strictly increasing and lie strictly between zero and one.

At every index $q$, at most one correction $f_j(q)$ is nonzero. Hence for
$S\ne\varnothing$, writing $s=|S|$,
\[
 \prod_{j\in S}(t_q+f_j(q))
   =t_q^s+t_q^{s-1}\sum_{j\in S}f_j(q).
\]
Averaging and using \eqref{eq:correct} gives
\[
 \frac1K\sum_q\prod_{j\in S}p_{j,q}
   =\mu_s(t)+s a_s(t)=\frac1{s+1}.
\]
The empty-set identity is automatic. Expanding the complementary factors
$1-p_{j,q}$ proves the claimed full comparison law.
\end{proof}

\begin{corollary}
The parameters in Theorem~\ref{thm:main} can be realized by actual finite
uniform dice: there are $m+1$ dice, one with $K=O(m^2)$ faces, for which
all comparisons with that distinguished die have exactly the same joint
law as under a uniform random ranking.
\end{corollary}
\begin{proof}
Give the distinguished die ordered faces $q_1<\cdots<q_K$.
For die $j$, assign its rational masses
\[
 p_{j,1},\ p_{j,2}-p_{j,1},\ldots,p_{j,K}-p_{j,K-1},\ 1-p_{j,K}
\]
to the $K+1$ gaps separated by those faces. Clearing denominators realizes
these masses by equally likely faces; choose distinct labels for all faces
inside each gap. Then $\Pr(D_j<q_q)=p_{j,q}$, as required.
This assertion concerns only comparisons with the distinguished die.
\end{proof}
\end{document}
